\chapter{Experimento en tiempo real con juegos}
\label{cap:experimentoTiempoReal}

\section{Participantes}
\label{sec:rt_participantes}

Al igual que en el \hyperref[cap:protocoloAdquisicion]{experimento de adquisición de grabaciones por trials}, todos los participantes respondieron el mismo cuestionario previo por medio de la aplicación. Además, los usuarios completaron un segundo cuestionario para medir sus niveles de satisfacción con el sistema, dicho formulario se pasó a cada usuario tantas veces como pruebas se realizaron con juegos.

Entre los participantes seleccionados se encuentran personas con experiencia previa en MI (gente que formó parte de pruebas anteriores de este mismo proyecto) al igual que otras sin experiencia. Los criterios de selección seguidos resultan los mismos que en el anterior experimento.

\section{Equipamiento y Configuración EEG}
\label{sec:rt_equipamiento}

Tanto el equipo EEG como la configuración de los electrodos se mantiene con respecto a los anteriores capítulos.

\section{Entorno experimental}
\label{sec:rt_entorno}

Se conservó tanto la sala como las condiciones del previo experimento.

\section{Paradigma Experimental}
\label{sec:rt_paradigma}

\subsection{Estructura del experimento}
\label{sec:rt_estructura}

El experimento sigue una organización muy concreta y ordenada, con estimaciones de duración para cada apartado:

\begin{itemize}
    \item \textbf{Explicación al usuario}
    \item \textbf{Grabación de \hyperref[cap:protocoloAdquisicion]{experimento visual}}
    \item \textbf{Descanso}
    \item \textbf{Entrenamiento \textit{Fine-tune}}
\end{itemize}

\subsection{Instrucciones al participante}
\label{subsec:rt_instrucciones}

\section{Recolección de métricas y resultados}
\label{sec:rt_metricas}
% Nota mental: Equivalente a "Formato de los datos" del Cap 5, pero enfocado en logs del sistema, tiempos de supervivencia en los juegos, uso del aimbot, etc.